\section{Dynamic Fit: Learn, Recruit, or Pivot}

\subsection{4 Responses to a Capability Gap}

The 1st mistake is thinking you must already be an ideal founder. The 2nd is thinking any gap can be overcome by unguided enthusiasm.

There are 4 basic responses to a capability gap:

\begin{enumerate}[label=\textbf{\arabic*.},leftmargin=11pt,itemsep=1.5pt,topsep=1.5pt,parsep=0pt]
  \item \textbf{Learn it.} Appropriate when the skill is learnable within a reasonable period, critical to founder judgment, and affordable to acquire. A technical founder should learn basic customer discovery and sales rather than assuming distribution can be outsourced before value is understood.
  \item \textbf{Hire or contract for it.} Appropriate when the requirement is well specified, the budget can support the cost, and external output can be audited. Bookkeeping, standard legal filings, design assets, or routine code integration fit this path.
  \item \textbf{Find a true partner or cofounder.} Appropriate when the missing capability is strategically central, requires ongoing creative judgment, and demands aligned long-term equity ownership. Cofounder ties carry high coordination and dissolution risks; equity should never be traded for vague promises of warm introductions.
  \item \textbf{Change the opportunity or business model.} \textbf{Changing the opportunity or business model is often the least romantic and most intelligent response to a major capability gap.} A founder with world-class technical teaching skill but modest capital should launch remote cohorts rather than attempting to lease a costly physical campus.
\end{enumerate}

Ask of each gap: \textit{How quickly can it be closed? What will it cost? How certain is the solution? Does closing it create a reusable capability?} If the gap is massive, expensive, and unrelated to your core strengths, the opportunity is poorly matched to your current starting position.

\subsection{T-Shaped Competence and Learning Velocity}

The goal is not omniscient genius across all functions. Early founders need baseline operational literacy: enough breadth to discover problems, speak to buyers, frame offers, deliver an initial result, and manage working capital. Beyond that floor, a distinctive spike creates advantage.

1 practical pattern is the \textbf{T-shaped competence profile}: deep specialized expertise in an important technical or domain function combined with practical literacy across sales, finance, product design, and hiring.

Being able to learn rapidly also matters. Still, \textbf{``I learn fast'' is a dangerous personality claim; treat learning velocity as an empirical track record to verify through concrete episodes and live trials.}
